\documentclass[10pt,a4paper]{amsart}
\usepackage[margin=19mm]{geometry}
\usepackage{amsmath,amssymb,mathtools}
\usepackage[scr=rsfs,cal=euler]{mathalfa}
\usepackage{xfrac}
\usepackage[unicode,hidelinks,bookmarks=false]{hyperref}
\newcommand{\ie}{i.e.\ }
\newcommand{\cf}{cf.\ }
\newcommand{\resp}{respectively\ }
\newcommand{\Subsection}{\subsection}
\providecommand{\nobreakdash}{\nobreak\hbox{-}}
\setlength{\emergencystretch}{2em}
\multlinegap0pt
\usepackage{amssymb}
\usepackage{mathtools}
\usepackage{xcolor}
\usepackage[shortlabels]{enumitem}
\usepackage{tikz}
\usepackage{float}
\newtheorem{proposition}{Proposition}
\newtheorem{remark}{Remark}
\newtheorem{lemma}{Lemma}
\usepackage{cleveref}
\crefname{proposition}{Proposition}{Propositions}
\crefname{remark}{Remark}{Remarks}
\crefname{lemma}{Lemma}{Lemmas}
\def\D{\Delta}
\def\F{\mathcal{F}}
\def\N{\mathcal{N}}
\def\l{\langle}
\def\r{\rangle}
\title{Square-free powers of Cohen-Macaulay simplicial forests}
\author{Kanoy Kumar Das, Amit Roy, Kamalesh Saha}
\date{Source: arXiv:2502.18396v2 (CC BY 4.0)}
\begin{document}
\maketitle

\noindent\emph{Source and licence.} Square-free powers of Cohen-Macaulay simplicial forests. Version arXiv:2502.18396v2 (CC BY 4.0), distributed by arXiv under the Creative Commons Attribution 4.0 International licence, http://creativecommons.org/licenses/by/4.0/, as recorded in the arXiv licence metadata for this exact version and shown on the abstract page https://arxiv.org/abs/2502.18396v2. Full licence notice: \texttt{licenses/arxiv-2502.18396-CC-BY-4.0.txt}.\par\smallskip

\noindent\textit{Editorial scope.} Counted unit: the article's lemma existence (source0.tex lines 362--364). The article's complete original proof of it is delivered with it (source0.tex lines 366--377, one \texttt{proof} environment of 12 lines). No mathematics is written, repaired or re-derived: the statement and the proof are the article's own lines, in the article's own notation, and no retained line is substituted or re-flowed.  Retained uncounted as the source-local context the proof needs: lines 238--240; lines 352--358.  Nothing was omitted from any retained span, so the package carries no omission record.  No retained line contains a citation, so the wrapper carries no bibliography and no bibliography entry.  No retained line contains a figure environment, an image-inclusion command or drawing code, so no image is delivered and the package declares no asset dependency.  Editorial formatting only, in addition: an AMS \texttt{amsart} preamble that restates the article's own declarations and macros used by the retained lines, namely the environments \texttt{proposition}, \texttt{remark}, \texttt{lemma} on separate counters, together with the standard packages those lines need.  The retained environments print in this document as Proposition 1, Remark 1, Lemma 1 in order of appearance, whereas the article prints the same items with the numbering of its own full document.  The complete native source of the article as distributed in the arXiv e-print for this exact version is bundled unchanged as \texttt{sources/173850/source0.tex}.
\begin{proposition}\label{goof leaf prop}
    Let $\D$ be a simplicial complex and $F\in \D$ be a leaf with $\N_{\D}(F)=\{F_1,\ldots , F_m\}$. Then $F$ is a good leaf of $\D$ if and only if there is a permutation $\{j_1,\ldots,j_m\}$ of the set $\{1,\ldots,m\}$ such that $F\cap F_{j_1}\supseteq F\cap F_{j_2}\supseteq\dots\supseteq F\cap F_{j_m}$.
\end{proposition}

    \begin{remark}\label{CM special leaf remark}
        Let $\Delta=\l F_1,\ldots,F_r\r\cup\l G_1,\ldots,G_s\r$ be a Cohen-Macaulay simplicial forest with the set of leaves $\{F_1,\ldots,F_r\}$. Since $F_i\cap F_j=\emptyset$, it is easy to see that a leaf $F$ of $\Delta$ is a special leaf if and only if the following implication holds:
        \[
G_l\cap G_m\neq\emptyset \ \Rightarrow \ (G_l\cap G_m)\setminus F\neq \emptyset,
        \]
    where $l,m\in[s]$ and $l\neq m$.
    \end{remark}

\begin{lemma}\label{special leaf existence}
    Let $\Delta$ be a Cohen-Macaulay simplicial forest. Then, $\Delta$ contains at least one special leaf.
\end{lemma}

\begin{proof}
    Let $\Delta=\l F_1,\ldots,F_r\r\cup\l G_1,\ldots,G_s\r$ be a Cohen-Macaulay simplicial forest, and set $\D'=\l G_1,\ldots,G_s\r$. Take any good leaf of $\D'$, say $G_1\in \F(\D')$. Without loss of generality, let $\N_{\D'}[G_1]=\{G_1,\ldots , G_t\}$ for some $t\leq s$. Since $G_1$ is a good leaf, by \Cref{goof leaf prop}, we can assume, without loss of generality, that $G_1\cap G_2\supseteq G_1\cap G_3\supseteq \cdots \supseteq G_1\cap G_t$. Now, take any leaf $F$ of $\D$ such that $F\cap G_1\neq \emptyset$. We first observe that $\N_{\D}(F)\subseteq \N_{\D'}[G_1]$. Indeed, if $F\cap G_j\neq \emptyset$ for some $j\in [s]\setminus\{1\}$, then since $F$ is a good leaf of $\D$, we have either $F\cap G_1\supseteq F\cap G_j\neq \emptyset$, or $F\cap G_j\supseteq F\cap G_1\neq \emptyset$. Thus in any case, $G_1\cap G_j\neq \emptyset$, which in turn, shows that $G_j\in \N_{\D'}[G_1]$.
    
    First, we consider the case when $F\nsupseteq G_1\cap G_t$. We claim that, in this case, $F$ is a special leaf. Note that, by \Cref{CM special leaf remark}, it is enough to show that if $G_l\cap G_m\neq \emptyset$ for some distinct $l,m\in [s]$, then $(G_l\cap G_m)\setminus F\neq \emptyset$. 
    
    \noindent \textbf{Case-I:} Either $G_l\notin \N_{\D'}[G_1]$, or $G_m\notin \N_{\D'}[G_1]$. Without loss of generality, assume $G_l\notin \N_{\D'}[G_1]$. Then we have $G_l\cap F=\emptyset$ since $\N_{\D}(F)\subseteq \N_{\D'}[G_1]$. Hence, $(G_l\cap G_m)\cap F=\emptyset$, which in turn, implies that $(G_l\cap G_m)\setminus F\neq \emptyset$, where $i\in [s]$.
    
     \noindent \textbf{Case-II:} $G_l,G_m\in \N_{\D'}[G_1]$. Then, since $G_1$ is a good leaf in $\D'$, without loss of generality, we can assume that $G_1\cap G_l\supseteq G_1\cap G_m\supseteq G_1\cap G_t$. Since $F\nsupseteq G_1\cap G_t$, there is some $x\in G_1\cap G_t$ such that $x\notin F$. Observe that $x\in G_l\cap G_m$, and $x\notin F$. This shows that $x\in (G_l\cap G_m)\setminus F$, where $i\in [s]$. This completes the proof of the claim.

     \noindent
     Finally, if $F\supseteq G_1\cap G_t$, then take any other leaf $F'\neq F$ of $\D$ such that $F'\cap G_1\neq \emptyset$. Such a leaf $F'$ exists since $G_1$ is a facet of $\D$. In this case, proceeding as before, we have $\N_{\D}(F')\subseteq \N_{\D'}[G_1]$. Now since $F\cap F'=\emptyset$, we must have $F'\nsupseteq G_1\cap G_t$. Then, we can proceed in a similar way as in the previous case to show that $F'$ is a special leaf of $\D$. Thus, $\D$ contains at least one special leaf.
\end{proof}

\end{document}
